\section{Comparative Synthesis}
\label{sec:comparative_synthesis}

\subsection{Structured Comparison}
To compare studies that use very different tasks and metrics, each study was analysed along the same dimensions: the role of the LLM, the task and data, the models, the language, the evaluation and the main finding. All numbers were checked against the original publications. Table~\ref{tab:comparison} shows the result. Two observations stand out. First, the simulation, opinion and generation studies work with English prompts and texts only (and the opinion study only with US groups), so their findings cannot be generalised to other languages or populations without further work. Second, each study uses its own evaluation (distribution comparison, human annotation or F1), so their results can be compared qualitatively but not on one numerical scale. % SOURCE: English-only: Aher (English prompts, Fig. 1 p.2); Santurkar US Pew surveys p.2; Vykopal "in the English language" p.1. (SemEval corpus language not stated explicitly in the task paper -> not claimed)

% Table 1: qualitative comparison of the four reviewed studies. Every entry is traceable to the
% source pages given in the section texts (Sections 3-6). Replaces the former F1 table (unverifiable values removed).
\begin{table}[t]
\centering
\footnotesize
\renewcommand{\arraystretch}{1.15}
\begin{tabular}{@{}p{1.9cm}p{3.0cm}p{2.7cm}p{1.0cm}p{4.9cm}@{}}
\toprule
\textbf{Study} & \textbf{Method / task} & \textbf{Data \& models} & \textbf{Lang.} & \textbf{Main finding} \\
\midrule
Aher et al.~\cite{aher2023using} (2023) & Turing Experiments: simulate participants of 4 human-subject studies & Own variants of classic studies; GPT-3 family, ChatGPT, GPT-4 & EN & 3 of 4 findings replicated by the largest models; \emph{hyper-accuracy distortion} in wisdom-of-crowds task \\
\addlinespace
Santurkar et al.~\cite{santurkar2023whose} (2023) & Compare LM and human opinion distributions ($W_1$-based alignment) & OpinionQA: 1,498 questions, 60 US groups; 9 LMs & EN (US) & Substantial misalignment; RLHF models shift towards liberal, educated, wealthy groups; steering gains modest \\
\addlinespace
Vykopal et al.~\cite{vykopal2024disinformation} (2024) & Prompt LLMs to write disinformation articles; human annotation; detector test & 20 narratives, 5 categories, 1,200 texts; 10 LLMs & EN & Several models write convincing disinformation; large differences in safety; detectors up to 0.82 macro-F1 \\
\addlinespace
Gaeta et al.~\cite{gaeta2025towards} (2025) & Prompted LLMs (role-play, few-shot, chain-of-thought) detect propaganda techniques and spans & SemEval-2020 Task~11 news articles; proprietary and open LLMs & n/s$^{*}$ & Reports notable improvements over state-of-the-art methods (scores not verified here$^{*}$) \\
\bottomrule
\end{tabular}
\caption{Qualitative comparison of the four reviewed studies (own analysis). $^{*}$Based on the publisher abstract only; the full text could not be accessed (see Section~\ref{sec:limitations}). n/s = not stated in the accessible text.}
\label{tab:comparison}
\end{table}


\subsection{Answers to the Research Questions}
\begin{description}
    \item[RQ1 (simulation fidelity).] Instruction-tuned models can reproduce the direction of several classic behavioural findings, and larger models do so more faithfully~\cite{aher2023using}. They deviate systematically where human answers are naturally spread out: in the wisdom-of-crowds task, recent aligned models give the exact answer with no variation at all (Section~\ref{sec:persona_simulation}). The own experiment shows that instruction tuning alone does not produce this distortion in a 0.5B model; the prompt format changed the spread, but not towards correct answers (Section~\ref{sec:experiment}).
    \item[RQ2 (opinion alignment).] LMs do not reflect the US population as a whole; human-feedback tuning shifts them towards liberal, educated and wealthier groups, and some groups (e.g.\ 65+, widowed) are poorly represented by all models. Prompting towards a group helps only modestly~\cite{santurkar2023whose}.
    \item[RQ3 (generation).] Several current LLMs write convincing disinformation articles on request, but models differ strongly, and open models are not necessarily less safe than proprietary ones~\cite{vykopal2024disinformation}. Existing detectors recognise the generated articles only under favourable conditions.
    \item[RQ4 (detection).] Prompted LLMs extend propaganda detection from trained classifiers to prompt-based analysis with chain-of-thought reasoning and span localisation, and are reported to improve on earlier methods~\cite{gaeta2025towards}. The size of the improvement could not be verified here, and the shared-task results show that the task remains difficult~\cite{dasanmartino2020semeval}.
\end{description}

\subsection{The Generation--Detection Asymmetry}
Taken together, the studies on generation and detection point to an asymmetry. On the generation side, Vykopal et al.\ show that some widely available models write convincing disinformation articles with almost no resistance~\cite{vykopal2024disinformation}, which supports earlier warnings that language models make such content cheaper to produce~\cite{buchanan2021truth}. On the detection side, the results are weaker and more conditional: the best machine-text detectors reach a macro-F1 of about 0.8 only under optimal thresholds (Figure~\ref{fig:detectors}), and in fine-grained propaganda analysis the best SemEval-2020 systems reached 51.55 F1 for span identification and 62.07 micro-F1 for technique classification~\cite{dasanmartino2020semeval}. % SOURCE: Vykopal Table 6 p.8; Da San Martino Table 5 p.14, Table 6 p.15
LLM-based detectors such as the system of Gaeta et al.\ report improvements on this benchmark~\cite{gaeta2025towards}, but they need several prompting steps for every document. In practice, a defender must check every text, while a producer of disinformation only needs some texts to go unnoticed. This asymmetry is an argument of this paper derived from the reviewed results, not a measured quantity.

\subsection{The Controllability Gap}
A second tension appears when the persona studies are compared with the detection study. Prompting is used successfully to direct LLMs towards an analytic task, namely finding propaganda techniques~\cite{gaeta2025towards}. When the same kind of prompting is used to make a model answer like a specific group of people, however, the effect is only modest and the representativeness problems remain~\cite{santurkar2023whose}. Aher et al.\ observe a related effect: the more aligned models are, the less they reproduce the natural spread of human answers~\cite{aher2023using}. This paper calls this contrast the \emph{controllability gap}: instructions change \emph{what task} a model carries out more easily than \emph{whose perspective} it takes. This is an interpretation across studies with different tasks and models, and it should be tested directly in future work.
